\chapter{Background \& Theory}
\label{chp:background}

\section{Visible Light Communication \& Sensing}

Visible light communication uses modulated light sources to transmit data, 
enabling a light source to serve simultaneously as illumination and a detectable 
navigation beacon.

\hlnotein{The spatial distribution of light intensity from a point source 
is described by the Lambertian model, which defines how intensity 
decays as a function of distance and angle from the source.}{Instead of dive into the Lambertian model,
I would instead just talk about it a bit, a more important point is that due to this Lambertian mode, 
how it affects or make other approach limited to multiple light overlapping to make it work.}

\hlnotein{Photodiodes convert incident light intensity into a measurable electrical signal, 
and their directional sensitivity makes them suitable for estimating the angle of incidence of a light source.}{I think this
paragraph can instad talk about what is the difference of our approach. What I mean is previous approach use light as
positioning system, our approach is more close to navigation system.}

\section{Bearing Angle Estimation}

\hlnotein{Prior work introduced a bearing angle estimation algorithm using a ring of 8 photodiodes 
that computes the relative angle toward a light source from the differential intensity readings across the array.}{No
need to talk about "8" phtodiodes here. more important thing is how to use sensor array to estimate the direction. 
Though we are almost the first one to do this with light, there are papers doing similar approach using other sensors.}

The bearing angle provides a single scalar direction toward the light source but 
contains no spatial information about the light environment; \hlnotein{when deploying multiple beacons,
their modulation frequencies must be chosen to avoid harmonic relationships, 
as the FFT cannot distinguish a frequency from its multiples.}{Too detail here, i would suggest to put it in the 
algorithm design. Here we are talking about more broader way of how to calcultae it. and the limitation.}

\section{Gradient-Based Source Seeking}

\hlnotein{Gradient-based source seeking navigates an agent toward a scalar field maximum by estimating and 
following the local gradient of that field.}{too compact, split to two sentences. this mean we will also spend a bit more
talking about this later.}

\hlnotein{Plane-fitting gradient estimation approximates the local gradient by fitting a least-squares plane 
to nearby light intensity measurements, yielding a continuous gradient direction without 
requiring deliberate perturbations.}{several terms is too detail to put here to explain how you make it work,
it should also be put later. at this point try to make it more easy for people to understand. something like
"by collecting the intensity overtime, a simple gradient estimation can be apply to estimate the direction of the light source.}

\hlnote{\section{Simulation Environment}}{I would put this more up, its ok if its short, but why we are using it? 
what is the advantage of using it? then talk about the good think about it.}

The simulation environment used in this thesis was developed in prior work and provides a physics-accurate Webots model of the 
drone platform with realistic light sensor simulation; for a full description the reader is referred to~\cite{}.
